\documentclass[10pt,a4paper]{amsart}
\usepackage[margin=19mm]{geometry}
\usepackage{amsmath,amssymb,mathtools}
\usepackage[scr=rsfs,cal=euler]{mathalfa}
\usepackage{xfrac}
\usepackage[unicode,hidelinks,bookmarks=false]{hyperref}
\newcommand{\ie}{i.e.\ }
\newcommand{\cf}{cf.\ }
\newcommand{\resp}{respectively\ }
\newcommand{\Subsection}{\subsection}
\providecommand{\nobreakdash}{\nobreak\hbox{-}}
\setlength{\emergencystretch}{2em}
\multlinegap0pt
\usepackage{amssymb}
\usepackage{float}
\usepackage{cancel}
\usepackage{xcolor}
\usepackage{mathtools}
\usepackage{bm}
\usepackage[shortlabels]{enumitem}
\newtheorem{proposition}{Proposition}
\newtheorem{theorem}{Theorem}
\newtheorem{lemma}{Lemma}
\title{Geodesics on Grushin spaces}
\author{Michael Albert, Samu\"el Borza, Maria Gordina}
\date{Source: arXiv:2509.03411v2 (CC BY 4.0)}
\begin{document}
\maketitle

\noindent\emph{Source and licence.} Geodesics on Grushin spaces. Version arXiv:2509.03411v2 (CC BY 4.0), distributed by arXiv under the Creative Commons Attribution 4.0 International licence, http://creativecommons.org/licenses/by/4.0/, as recorded in the arXiv licence metadata for this exact version and shown on the abstract page https://arxiv.org/abs/2509.03411v2. Full licence notice: \texttt{licenses/arxiv-2509.03411-CC-BY-4.0.txt}.\par\smallskip

\noindent\textit{Editorial scope.} Counted unit: the article's lemma 3D (source0.tex lines 1355--1374). The article's complete original proof of it is delivered with it (source0.tex lines 1376--1560, one \texttt{proof} environment of 185 lines). No mathematics is written, repaired or re-derived: the statement and the proof are the article's own lines, in the article's own notation, and no retained line is substituted or re-flowed.  Retained uncounted as the source-local context the proof needs: lines 1095--1105; lines 1107--1122; lines 1145--1162; lines 1300--1328.  Nothing was omitted from any retained span, so the package carries no omission record.  No retained line contains a citation, so the wrapper carries no bibliography and no bibliography entry.  No retained line contains a figure environment, an image-inclusion command or drawing code, so no image is delivered and the package declares no asset dependency.  Editorial formatting only, in addition: an AMS \texttt{amsart} preamble that restates the article's own declarations and macros used by the retained lines, namely the environments \texttt{proposition}, \texttt{theorem}, \texttt{lemma} on separate counters, together with the standard packages those lines need.  The retained environments print in this document as Proposition 1, Theorem 1, Lemma 1 in order of appearance, whereas the article prints the same items with the numbering of its own full document.  The complete native source of the article as distributed in the arXiv e-print for this exact version is bundled unchanged as \texttt{sources/184410/source0.tex}.
\begin{equation}\label{Conjectured Cut}
    \tau_j(\phi)
    :=
    \min
    \left\{
    t>0:
    |\omega_j|\int_0^t\xi_j^2(s)\,ds
    =
    \pi_{\alpha_j}
    \right\}.
\end{equation}

\begin{proposition}\label{tan function solution}
    Let \(\alpha\geq1\) and
    \[
        \phi\in[0,2\pi_\alpha)
        \setminus
        \left\{
        \frac{\pi_\alpha}{2},
        \frac{3\pi_\alpha}{2}
        \right\}.
    \]
    Then the equation
    \begin{equation}\label{tan equation}
        \tan_\alpha(\phi+\zeta)-\tan_\alpha(\phi)=\zeta
    \end{equation}
    has no nonzero solution with \(|\zeta|<\pi_\alpha\).
\end{proposition}

\begin{theorem}\label{Determinant}
    Let \(x^0\in\mathbb G_\alpha^{n+1}\) be Riemannian, and let
    \(\phi\in\mathcal R_\alpha^0\) correspond to a non-trivial geodesic. Define
    \[
        \tau(\phi):=\min_{j\in J}\tau_j(\phi),
    \]
    with \(\tau_j\) as in \eqref{Conjectured Cut}. Then
    \begin{equation}\label{conjugate inequality}
        0<\tau(\phi)\leq t_{\operatorname{con}}(\phi).
    \end{equation}
    Moreover, equality holds if and only if there is \(j\in J\) such that
    \(\tau=\tau_j\) and
    \[
        \phi_j=\frac{\pi_{\alpha_j}}{2}
        \quad\text{or}\quad
        \phi_j=\frac{3\pi_{\alpha_j}}{2}.
    \]
\end{theorem}

\begin{equation}\label{det formula 3D}
    \begin{aligned}
        D(t,\phi_1,\phi_2)
         & =
        -
        \frac{|x_0|^{\alpha+1}|y_0|^{\beta+1}}
        {|\sin_\alpha(\phi_1)|^{\alpha+1}|\sin_\beta(\phi_2)|^{\beta+1}}
        \\
         & \quad\times
        \left[
            \cos_\alpha(A)
            \left(\cot_\alpha(\phi_1)\omega_1t+1\right)
            -
            \cot_\alpha(\phi_1)\sin_\alpha(A)
            \right]
        \\
         & \quad\times
        \left[
            \cos_\beta(B)
            \left(
            \cot_\beta(\phi_2)\omega_2
            \int_0^t x^{2\alpha}(s)\,ds
            +1
            \right)
            -
            \cot_\beta(\phi_2)\sin_\beta(B)
            \right].
    \end{aligned}
\end{equation}

\begin{lemma}\label{Con Times 3D}
    Let \(\gamma(t)=\exp_{q_0}(t\lambda_0)\) be a geodesic in
    \(\mathbb G^3_{(\alpha,\beta)}\) starting from a Riemannian point. Then \(\gamma\) has no
    conjugate time in \((0,\tau)\). Moreover, \(t_{\operatorname{con}}=\tau\) if either
    \[
        \tau=\tau_1
        \quad\text{and}\quad
        \phi_1=\frac{\pi_\alpha}{2},
    \]
    or
    \[
        \tau=\tau_2
        \quad\text{and}\quad
        \phi_2\in
        \left\{
        \frac{\pi_\beta}{2},
        \frac{3\pi_\beta}{2}
        \right\}.
    \]
\end{lemma}

\begin{proof}
    For \(w_0\neq0\), the claim follows from Theorem~\ref{Determinant}. We treat the boundary
    cases \(w_0=0\).

    First assume \(w_0=0\) and \((u_0,v_0)\neq(\pm1,0)\). In spherical coordinates this
    corresponds to \(\phi_2\to0\) or \(\phi_2\to\pi_\beta\), with
    \(\phi_1\notin\{0,\pi_\alpha\}\). We consider the limit as \(\phi_2\to0\); the case
    \(\phi_2\to\pi_\beta\) is identical. Let
    \[
        I(t):=\int_0^t x^{2\alpha}(s)\,ds,
        \qquad
        \varepsilon:=\frac{\delta_1}{y_0}.
    \]
    Then
    \[
        \omega_2=\varepsilon\sin_\beta(\phi_2),
        \qquad
        B=\phi_2+\varepsilon\sin_\beta(\phi_2)I(t).
    \]
    Using the Taylor expansions of \(\sin_\beta\) and \(\cos_\beta\) at \(0\), we have
    \[
        \sin_\beta(B)
        \sim
        \sin_\beta(\phi_2)\bigl(1+\varepsilon I(t)\bigr),
        \qquad
        \cos_\beta(B)\sim 1,
        \qquad
        \phi_2\to0.
    \]
    The \(\phi_2\)-factor in \eqref{det formula 3D} has the indeterminate form
    \(0/|\sin_\beta(\phi_2)|^{\beta+1}\). Differentiating numerator and denominator gives
    \[
        \begin{aligned}
             & \frac{d}{d\phi_2}
            \left[
                \cos_\beta(B)
                \left(\cot_\beta(\phi_2)\omega_2 I(t)+1\right)
                -
                \cot_\beta(\phi_2)\sin_\beta(B)
                \right]
            \\
             & \qquad\sim
            -\beta\sin_\beta^{2\beta-1}(\phi_2)
            \left[
                (1+\varepsilon I(t))^3
                +
                \varepsilon I(t)
                -
                (1+\varepsilon I(t))
                \right],
        \end{aligned}
    \]
    while
    \[
        \frac{d}{d\phi_2}
        |\sin_\beta(\phi_2)|^{\beta+1}
        \sim
        (\beta+1)|\sin_\beta(\phi_2)|^\beta .
    \]
    Thus the spherical determinant degenerates by a factor
    \(|\sin_\beta(\phi_2)|^{\beta-1}\). This is exactly the degeneracy of the spherical
    coordinate system. Indeed,
    \[
        \left|
        \frac{\partial(u_0,v_0,w_0)}
        {\partial(t,\phi_1,\phi_2)}
        \right|
        =
        \frac{
            t^2\alpha\beta
            \sin_\alpha^{2\alpha-1}(\phi_1)
            |\sin_\beta(\phi_2)|^{\beta-1}
        }{
            x_0^{2\alpha}|y_0|^\beta
        }.
    \]
    Therefore the Cartesian determinant
    \[
        D^{\operatorname{Cart}}\exp_{q_0}
        =
        \left|
        \frac{\partial(u_0,v_0,w_0)}
        {\partial(t,\phi_1,\phi_2)}
        \right|^{-1}
        D(t,\phi_1,\phi_2)
    \]
    has a finite limit as \(\phi_2\to0\). More precisely,
    \begin{equation}\label{e.cart-limit-w0}
        \begin{aligned}
            \lim_{\phi_2\to0}
            D^{\operatorname{Cart}}\exp_{q_0}
             & =
            C(\phi_1,t)
            \left[
                \cos_\alpha(A)
                \left(\cot_\alpha(\phi_1)\omega_1t+1\right)
                -
                \cot_\alpha(\phi_1)\sin_\alpha(A)
                \right]
            \\
             & \quad\times
            \left[
                (1+\varepsilon I(t))^3
                +
                \varepsilon I(t)
                -
                (1+\varepsilon I(t))
                \right],
        \end{aligned}
    \end{equation}
    where \(C(\phi_1,t)\neq0\) for \(t>0\) and
    \(\phi_1\notin\{0,\pi_\alpha\}\). The last bracket is
    \[
        (1+\varepsilon I)^3+\varepsilon I-(1+\varepsilon I)
        =
        \varepsilon I\bigl((1+\varepsilon I)^2+(1+\varepsilon I)+2\bigr).
    \]
    Since \(I(t)>0\) for \(t>0\), and \(\varepsilon\neq0\), this factor has no positive zero.
    Consequently, the only possible zeros of the limiting Cartesian determinant are the zeros
    of
    \[
        \cos_\alpha(A)
        \left(\cot_\alpha(\phi_1)\omega_1t+1\right)
        -
        \cot_\alpha(\phi_1)\sin_\alpha(A).
    \]
    By Proposition~\ref{tan function solution}, this factor does not vanish for
    \(0<t<\tau_1\). Since \(\tau_2=+\infty\) when \(w_0=0\), there are no conjugate times in
    \((0,\tau)\) for \(w_0=0\), \((u_0,v_0)\neq(\pm1,0)\).

    It remains to consider the straight-line geodesics \(\lambda_0=(\pm1,0,0)\). Write
    \(\eta=\pm1\). Then
    \[
        x(t)=x_0+\eta t,
        \qquad
        y(t)=y_0,
        \qquad
        z(t)=z_0.
    \]
    Use local Cartesian coordinates \((t,v_0,w_0)\) on the energy level near
    \((\eta,0,0)\), with \(u_0\) determined by
    \[
        u_0^2+x_0^{2\alpha}v_0^2+x_0^{2\alpha}y_0^{2\beta}w_0^2=1.
    \]
    At \(v_0=w_0=0\), one has
    \[
        \frac{\partial x}{\partial v_0}
        =
        \frac{\partial x}{\partial w_0}
        =
        0,
    \]
    \[
        \frac{\partial y}{\partial v_0}
        =
        \int_0^t (x_0+\eta s)^{2\alpha}\,ds,
        \qquad
        \frac{\partial y}{\partial w_0}=0,
    \]
    and
    \[
        \frac{\partial z}{\partial w_0}
        =
        y_0^{2\beta}
        \int_0^t (x_0+\eta s)^{2\alpha}\,ds,
        \qquad
        \frac{\partial z}{\partial v_0}=0.
    \]
    Since \(\partial_t x=\eta\), we obtain
    \begin{equation}\label{e.straight-line-det}
        \det
        \frac{\partial(x,y,z)}
        {\partial(t,v_0,w_0)}
        =
        \eta\,y_0^{2\beta}
        \left(
        \int_0^t (x_0+\eta s)^{2\alpha}\,ds
        \right)^2.
    \end{equation}
    This is nonzero for every \(t>0\). Hence straight-line geodesics have no conjugate
    points, which agrees with \(\tau_1=\tau_2=+\infty\).

    The equality statement follows from Theorem~\ref{Determinant} in the non-boundary cases
    and from the preceding limiting computation when \(w_0=0\).
\end{proof}

\end{document}
